\section{Measurement uncertainty: derivations}
\label{sec:annex-p15-measurement-uncertainty}

\begin{annexitem}{annex:p15.measurement.sample-covariance}{thm:p15-measurement-sample-covariance}{Covariance of a mean and its sample estimator}
\begin{proof}
Write \(W_k=Q_k-\mu\), \(\bar W=n^{-1}\sum_kW_k\).
Linearity gives \(E\bar Q=\mu\). Independence and centering give
\(E(W_kW_l^{\mathsf T})=0\) for \(k\ne l\); hence
\(E(\bar W\bar W^{\mathsf T})=C/n\).
Direct expansion, using \(\sum_kW_k=n\bar W\), gives
\[
 \sum_k(Q_k-\bar Q)(Q_k-\bar Q)^{\mathsf T}
 =\sum_k W_kW_k^{\mathsf T}-n\bar W\bar W^{\mathsf T}.
\]
Its expectation is \(nC-C=(n-1)C\), proving \(ES=C\).
Dividing by \(n\) gives the asserted covariance estimator. The displayed
scalar formula is its diagonal entry. Without independence, expand the
covariance of \(n^{-1}\sum Q_k\) bilinearly and retain all \(n^2\) terms.
Finite second moments justify every expectation used.
\end{proof}
\end{annexitem}

\begin{annexitem}{annex:p15.measurement.bounded-distributions}{thm:p15-measurement-bounded-distributions}{Uniform, triangular, trapezoidal and sinusoidal models}
\begin{proof}
Center at zero. Symmetry gives zero means. For the uniform model,
\[
 E T^2=\frac1{2a}\int_{-a}^a t^2\,dt=\frac{a^2}{3}.
\]
For the triangular model,
\[
 E T^2=\frac2{a^2}\int_0^a t^2(a-t)\,dt=\frac{a^2}{6}.
\]
Both densities integrate to one by the same elementary integrals.

Put \(a_1=a(1+\beta)/2\), \(a_2=a(1-\beta)/2\). For \(a_2>0\),
the convolution of independent uniforms on \([-a_1,a_1]\) and
\([-a_2,a_2]\) has density equal to the overlap length of these intervals
after translation, divided by \(4a_1a_2\). Its support is \([-a,a]\);
the overlap is constant for \(|t|\leq a_1-a_2=a\beta\) and decreases
linearly on either side. It is precisely the stated trapezoid.
Independence gives variance
\[
 (a_1^2+a_2^2)/3=a^2(1+\beta^2)/6.
\]
At \(\beta=1\), let the second summand be identically zero; at
\(\beta=0\) the preceding formula gives the triangular case.

For uniform phase \(\Phi\), integration on \([0,2\pi]\) gives
\(E\sin\Phi=0\) and \(E\sin^2\Phi=1/2\). Thus
\(\operatorname{Var}(c+a\sin\Phi)=a^2/2\).
On the interior of \([-a,a]\), the two inverse branches per full period
give the centered density \(1/(\pi\sqrt{a^2-t^2})\).
This also explains the source's U-shaped cyclic-temperature model.
\end{proof}
\end{annexitem}

\begin{annexitem}{annex:p15.measurement.asymmetric-model}{thm:p15-measurement-asymmetric-model}{A bounded entropy model with a fixed mean}
\begin{proof}
Set \(T=t-m\), so the interval is \([-a,b]\) with \(a=b_->0\),
\(b=b_+>0\). The normalizing integral
\(Z(\lambda)=\int_{-a}^b e^{-\lambda T}\,dT\) is positive and smooth
for every real \(\lambda\). Differentiation under an integral over a
compact interval is justified by uniform bounds on each derivative on a
compact \(\lambda\)-interval. The mean
\[
 h(\lambda)=-\frac{Z'(\lambda)}{Z(\lambda)}
\]
satisfies \(h'(\lambda)=-\operatorname{Var}_{q_\lambda}(T)<0\):
the density is positive on a nondegenerate interval.

As \(\lambda\to+\infty\), mass concentrates at \(-a\).
Indeed, for \(0<\varepsilon<a+b\), compare the integral on
\([-a+\varepsilon,b]\) with that on
\([-a,-a+\varepsilon/2]\). Their ratio is at most
\(2(a+b)e^{-\lambda\varepsilon/2}/\varepsilon\).
Boundedness of \(T\) then implies \(h(\lambda)\to-a\).
The corresponding reversed comparison gives \(h(\lambda)\to b\)
as \(\lambda\to-\infty\). Continuity and strict decrease give a unique
zero. At \(\lambda=0\), the mean is \((b-a)/2\); hence zero occurs there
exactly in the symmetric case.

To establish entropy optimality, let \(p\) be any density on the same
interval with mean zero. The function \(r\log r-r+1\) is nonnegative
for \(r\geq0\), with equality only at \(r=1\); this follows by
differentiation and continuity at zero. Integrating with \(r=p/q_\lambda\)
gives
\[
 D(p\Vert q_\lambda)=\int p\log(p/q_\lambda)\,dt\geq0.
\]
The integral is interpreted as an extended nonnegative divergence.
Since \(\log q_\lambda=-\lambda T-\log Z\) is bounded on the interval,
the common mean constraint implies
\(\int p\log q_\lambda=\int q_\lambda\log q_\lambda=-\log Z\).
Consequently the differential entropy of \(p\) is at most that of
\(q_\lambda\), including the case of entropy \(-\infty\).
Equality requires \(p=q_\lambda\) almost everywhere.

For \(\lambda\ne0\), integration of
\((T e^{-\lambda T})'\), using \(ET=0\), gives
\[
 b e^{-\lambda b}+a e^{\lambda a}=Z,\qquad
 \lambda Z=e^{\lambda a}-e^{-\lambda b}.
\]
Integration of \((T^2e^{-\lambda T})'\) gives
\[
 \operatorname{Var}(T)=
 \frac{a^2e^{\lambda a}-b^2e^{-\lambda b}}{\lambda Z}.
\]
Substituting the preceding two identities reduces this to
\(ab-(b-a)/\lambda\). At \(\lambda=0\), the uniform calculation
in \cref{thm:p15-measurement-bounded-distributions} applies separately.
\end{proof}
\end{annexitem}

\begin{annexitem}{annex:p15.measurement.affine-covariance}{thm:p15-measurement-affine-covariance}{Affine uncertainty propagation and valid covariance matrices}
\begin{proof}
Linearity of expectation gives the mean. Subtract it to obtain
\(Y-EY=J(X-m)\). Taking the outer product and expectation gives
\[
 E[(Y-EY)(Y-EY)^{\mathsf T}]=J\Sigma J^{\mathsf T}.
\]
Expanding its scalar entry gives the two sums in the statement.
Covariances of uncorrelated inputs vanish by definition, so independence
is not needed for that reduction.

For any real vector \(v\),
\(v^{\mathsf T}\Sigma v=E\{[v^{\mathsf T}(X-m)]^2\}\geq0\).
For any symmetric positive semidefinite assessed matrix \(U\), the same
quadratic-form argument gives
\(w^{\mathsf T}JUJ^{\mathsf T}w\geq0\). Restricting the quadratic form
of \(U\) to coordinates \(i,j\) gives
\(u_{ij}^2\leq u_{ii}u_{jj}\), including the fact that a zero diagonal
entry forces the corresponding row to vanish. Thus defined correlation
coefficients lie in \([-1,1]\).

If all relevant correlations equal \(+1\), then
\(u_{ij}=u_i u_j\), so
\(c^{\mathsf T}Uc=(\sum_i c_i u_i)^2\); the nonnegative square root
is the absolute value stated in the body. For zero off-diagonal entries,
the root-sum-of-squares expression follows instead.
\end{proof}
\end{annexitem}

\begin{annexitem}{annex:p15.measurement.common-correction}{thm:p15-measurement-common-correction}{Shared offset and scale under repeated observations}
\begin{proof}
By independence \(E(BQ)=bq\) and
\(E[(BQ)^2]=E(B^2)E(Q^2)=(v_B+b^2)(v_Q+q^2)\).
Subtracting \(b^2q^2\) gives
\[
 \operatorname{Var}(BQ)=q^2v_B+b^2v_Q+v_Bv_Q.
\]
Independence of \(A\) removes its covariance with \(BQ\), adding \(v_A\).
The formula for the mean follows by linearity.
For an average of independent identically distributed readings,
\cref{thm:p15-measurement-sample-covariance} gives \(v_Q=\sigma^2/n\).
The shared-offset difference and sum example follows directly by cancelling
\(C\) algebraically in the difference, retaining \(2C\) in the sum, and
using covariance additivity. It does not require normal distributions.
\end{proof}
\end{annexitem}

\begin{annexitem}{annex:p15.measurement.taylor-bound}{thm:p15-measurement-taylor-bound}{A remainder bound for the output mean and variance}
\begin{proof}
For a fixed allowed displacement \(\delta\), apply the fundamental theorem
of calculus twice to \(h(t)=f(m+t\delta)\) on \([0,1]\):
\[
 f(m+\delta)=f(m)+Df(m)\delta+
  \int_0^1(1-t)\delta^{\mathsf T}D^2f(m+t\delta)\delta\,dt.
\]
The Hessian bound gives \(|R|\leq M\|\delta\|^2/2\).
Taking expectations and using \(E(X-m)=0\) gives the asserted mean bound;
squaring the remainder bound and taking expectations gives \(\|R\|_2\leq r\).

All variables are square integrable, and
\[
 \operatorname{Var}(L+R)-\operatorname{Var}(L)
 =2\operatorname{Cov}(L,R)+\operatorname{Var}(R).
\]
Cauchy--Schwarz bounds the absolute covariance by
\(\sqrt{\operatorname{Var}(L)\operatorname{Var}(R)}\), while
\(\operatorname{Var}(R)\leq ER^2\leq r^2\).
This proves the variance bound.

If \(f\) is quadratic, its Hessian \(H\) is constant and the integral
remainder is exactly \(\delta^{\mathsf T}H\delta/2\). Taking expectations
gives the body's identity
\(Ef(X)-f(m)=\sum_{i,j}H_{ij}\Sigma_{ij}/2\), with only second moments
needed in that special case.
\end{proof}
\end{annexitem}

\begin{annexitem}{annex:p15.measurement.gaussian-expansion}{thm:p15-measurement-gaussian-expansion}{Quadratic and cubic contributions to Gaussian propagated variance}
\begin{proof}
For a centered normal \(Z\) of positive variance \(\sigma^2\),
symmetry gives zero odd moments. Integration by parts in its Gaussian
density gives
\[
 E Z^{r+2}=(r+1)\sigma^2 E Z^r \qquad(r=0,2,4).
\]
The boundary terms vanish because a fixed polynomial times the Gaussian
density tends to zero. Thus \(EZ^2=\sigma^2\),
\(EZ^4=3\sigma^4\) and \(EZ^6=15\sigma^6\).
If \(\sigma=0\), \(Z=0\) almost surely and these statements also hold.

Independence and these one-dimensional identities imply
\[
 E(Z_iZ_jZ_kZ_l)=
 \delta_{ij}\delta_{kl}\sigma_i^2\sigma_k^2+
 \delta_{ik}\delta_{jl}\sigma_i^2\sigma_j^2+
 \delta_{il}\delta_{jk}\sigma_i^2\sigma_j^2,
\]
where \(\delta_{ij}\) is one for equal indices and zero otherwise.
To verify it, if an index appears an odd number of times, both sides
vanish; two distinct pairs give the product of variances; four equal
indices give \(3\sigma_i^4\). These exhaust the cases.

It follows that \(EA=EC=0\) and
\(EB=\frac12\sum_iH_{ii}\sigma_i^2\).
The variance of \(A\) is \(\sum_i g_i^2\sigma_i^2\).
Expanding \(EB^2-(EB)^2\) with the fourth-moment identity cancels
the pairing internal to the two quadratic factors and leaves
\[
 \operatorname{Var}(B)=\frac12\sum_{i,j}H_{ij}^2\sigma_i^2\sigma_j^2.
\]
Similarly, symmetry of \(T\) and the three fourth-moment pairings give
\[
 \operatorname{Cov}(A,C)=
 \frac12\sum_{i,j}g_iT_{ijj}\sigma_i^2\sigma_j^2.
\]
The joint law of \(Z\) is invariant under \(Z\mapsto-Z\).
Therefore \(\operatorname{Cov}(A,B)=0\) and
\(\operatorname{Cov}(B,C)=0\); the relevant mixed products have odd
total degree. Expanding the variance of
\(\varepsilon A+\varepsilon^2B+\varepsilon^3C\) gives the displayed
identity. All sixth moments are finite, so \(\operatorname{Var}(C)\)
is finite. No omitted Taylor term has been declared zero for a
nonpolynomial function.
\end{proof}
\end{annexitem}

\begin{annexitem}{annex:p15.measurement.coverage-models}{thm:p15-measurement-coverage-models}{Coverage under different output distributions}
\begin{proof}
Standardizing a centered normal output gives
\(P(-k\leq Z\leq k)=\Phi(k)-\Phi(-k)=2\Phi(k)-1\).
For a uniform law, the interval of radius \(k a/\sqrt3\) occupies the
fraction \(k/\sqrt3\) of its support until that fraction reaches one.
For the triangular density and \(0\leq z\leq1\),
\[
 P(|T-c|\leq az)=\frac2{a^2}\int_0^{az}(a-t)\,dt=2z-z^2.
\]
Past the support boundary the probability is one.
Finally
\[
 \sigma^2=E[(Y-EY)^2]
 \geq k^2\sigma^2 P(|Y-EY|\geq k\sigma).
\]
Divide by \(k^2\sigma^2\) and take complements.
This proof uses the true mean and variance of the stipulated output law.
Replacing them with estimates is an additional statistical step.
\end{proof}
\end{annexitem}

\begin{annexitem}{annex:p15.measurement.effective-df}{thm:p15-measurement-effective-df}{Moment matching behind the Welch--Satterthwaite expression}
\begin{proof}
Independence gives \(EV=v\) and
\(\operatorname{Var}(V)=2\sum_i v_i^2/\nu_i\).
A chi-square variable of parameter \(\nu>0\) is a gamma variable of
shape \(\nu/2\) and rate \(1/2\); its mean and variance are
\(\nu\) and \(2\nu\), by
\cref{thm:p15-gamma-normalization}.
Hence \(v\chi_\nu^2/\nu\) has mean \(v\) and variance \(2v^2/\nu\).
Matching variances yields the displayed formula.
A general weighted sum of independent chi-square variables is not thereby
shown to have a chi-square law. In particular no Student pivot or exact
coverage follows from these two moments alone.
\end{proof}
\end{annexitem}

\begin{annexitem}{annex:p15.measurement.gauge-products}{thm:p15-measurement-gauge-products}{Thermal products in the end-gauge model}
\begin{proof}
Independence and \(EA=ED=0\) give \(E(A\Theta)=E(BD)=0\).
Moreover
\[
 E[(A\Theta)^2]=v_A(v_\Theta+\theta^2),\qquad
 E[(BD)^2]=v_D(v_B+b^2).
\]
The two products depend on disjoint independent pairs, so their covariance
is zero; \(E_0\) is independent of both. Adding their variances proves
the formula and adding their means gives \(EG=g_0\).

For the numerical case, in square nanometres the independent length and
comparison contribution is
\[
 v_0=25^2+\frac{13^2}{5}+\left(\frac{10}{2.57}\right)^2+
                 \left(\frac{20}{3}\right)^2
       \approx718.384719.
\]
The two first-order thermal terms are \(8.333333\) and \(275.520833\).
The additional product terms are
\[
 (0.05)^2\frac{10^{-12}}3(0.2^2+0.5^2/2)\,10^{18}=137.5,
\]
\[
 (0.05)^2\frac{4\,10^{-12}}3\frac{0.05^2}3\,10^{18}
 =2.7777778.
\]
The factor \(10^{18}\) converts square metres to square nanometres.
Their sum with the linear budget is approximately \(1142.516664\),
whose square root is \(33.801134\) nanometres.
The independent-product model gives an exact variance for the displayed
polynomial surrogate, while approximation of the physical rational model
remains a separate issue.

For completeness, if the multiplicative reference length is an independent
random \(S\), mean \(L\) and variance \(v_S\), and
\(F=A\Theta+BD\) is centered and independent of \(S\), then
\(\operatorname{Var}(SF)=(L^2+v_S)\operatorname{Var}(F)\) and
\(\operatorname{Cov}(S,SF)=0\).
Thus using \(L\) in place of the uncertain multiplier omits
\(v_S\operatorname{Var}(F)\). This is an additional product contribution,
not part of the first-order standard-length term.
The algebraic remainder of the rational length model in the body follows
by bringing its exact and approximate expressions to the common
denominator \(1+\alpha\theta\).
\end{proof}
\end{annexitem}

\begin{annexitem}{annex:p15.measurement.impedance-jacobian}{thm:p15-measurement-impedance-jacobian}{The electrical-output covariance calculation}
\begin{proof}
Differentiate the three coordinate functions with \(I>0\).
The first two phase derivatives are respectively
\(-(V/I)\sin\phi=-X_{\rm r}\) and
\((V/I)\cos\phi=R\); the third is zero.
The voltage and current derivatives give the other two columns of \(J\).
Applying \cref{thm:p15-measurement-affine-covariance} to the linear
surrogate gives \(JUJ^{\mathsf T}\).
For its third row the variance is
\[
 u_V^2/I^2+V^2u_I^2/I^4-2Vu_{VI}/I^3.
\]
Division by \(Z^2=V^2/I^2\), for \(V>0\), gives the relative expression.
The trigonometric identity gives \(R^2+X_{\rm r}^2=Z^2\).

For the five simultaneous rows, the covariance matrix of the input mean,
using volts, amperes and radians as the coordinates, is
\[
 U=\begin{pmatrix}
 1.0300\,10^{-5}&-1.0800\,10^{-8}&2.0700\,10^{-6}\\
 -1.0800\,10^{-8}&8.9700\,10^{-11}&-4.5950\,10^{-9}\\
 2.0700\,10^{-6}&-4.5950\,10^{-9}&5.6560\,10^{-7}
 \end{pmatrix}.
\]
Each entry is the finite sum in
\cref{thm:p15-measurement-sample-covariance}, with denominator \(5\cdot4\).
At the input means,
\[
 J\approx\begin{pmatrix}
 25.551544&-6496.728037&-219.846512\\
 43.978098&-11181.858090&127.732170\\
 50.862113&-12932.185644&0
 \end{pmatrix}.
\]
Multiplication gives the first-order output covariance in square ohms:
\[
 JUJ^{\mathsf T}\approx
 \begin{pmatrix}
 0.005051145&-0.012361383&-0.008150774\\
 -0.012361383&0.087368528&0.069333519\\
 -0.008150774&0.069333519&0.055854766
 \end{pmatrix}.
\]
Square roots of its diagonal and normalization of its off-diagonal entries
give the body table. Replacing \(U\) by its diagonal gives the last row.
Transforming each of the five original triples first and then applying
the sample-covariance formula gives the separate paired-output row.
The computation retains the observed pairing; it supplies no justification
for assigning that pairing to sequential, unrelated measurements.
\end{proof}
\end{annexitem}

\begin{annexitem}{annex:p15.measurement.transform-average}{thm:p15-measurement-transform-average}{Finite-sample transformation and averaging}
\begin{proof}
For affine \(f(q)=a+c^{\mathsf T}q\), linearity of a finite sum gives
the result. For quadratic \(f\), write \(\delta_k=q_k-\bar q\).
Its exact expansion is
\[
 f(q_k)=f(\bar q)+Df(\bar q)\delta_k+
                   \tfrac12\delta_k^{\mathsf T}H\delta_k.
\]
Averaging eliminates the linear term because \(\sum_k\delta_k=0\).
By the definition of \(U_{\rm mean}\),
\[
 \frac1n\sum_k(\delta_k)_i(\delta_k)_j
       =(n-1)(U_{\rm mean})_{ij}.
\]
Substitution gives the assertion.
For twice continuously differentiable \(f\), the proof of
\cref{thm:p15-measurement-taylor-bound} gives instead the exact difference
\[
 \frac1n\sum_k\int_0^1(1-t)\,
  \delta_k^{\mathsf T}D^2f(\bar q+t\delta_k)\delta_k\,dt.
\]
A probability expectation of a transformed random input uses its own
covariance, rather than the differently normalized empirical covariance
of a sample mean. This proves the distinction used in the body.
\end{proof}
\end{annexitem}
